% !TeX root = ../main.tex

\xchapter{结论和展望}{Conclusion and Future Work}

\xsection{结论}{Conclusion}

本文围绕复杂室内环境下的非视觉人体感知问题展开研究，重点关注人体活动识别与人体姿态重构与预测两个任务。前者面向离散活动语义理解，后者面向连续人体运动状态建模，由宏观活动类别识别延伸至细粒度姿态时序建模，共同服务于隐私友好、部署灵活且适用于复杂室内环境的人体感知需求。

在人体活动识别任务中，针对单一非视觉模态信息表达有限、IMU 与 WiFi CSI 异构模态难以有效融合以及跨受试者泛化能力不足等问题，本文提出了基于多模态对齐与融合的人体活动识别方法 WIHAR。该方法通过双分支 Mamba 时序编码器分别提取 IMU 与 WiFi CSI 特征，并结合跨模态双向对比学习与交叉注意力机制实现模态对齐和特征融合。实验结果表明，在 XRFV2 数据集上，WIHAR 的活动识别准确率达到 95.81\%，宏平均 F1、宏平均精确率和宏平均召回率分别达到 95.29\%、95.48\% 和 95.13\%，均优于多种代表性基线方法；在 OctoNet 数据集上的补充验证结果也表明，WIHAR 在其他包含 WiFi 和 IMU 双模态数据集上仍保持较好的识别性能。此外，基于类别原型的跨受试者适配策略能够在少量目标受试者标注样本条件下提升模型对未见个体的适应能力，使跨受试者设置下的平均识别准确率由 72.57\% 提升至 90.99\%。

在人体姿态重构与预测任务中，针对现有基于 IMU 的人体姿态方法多关注当前姿态重构、缺少未来姿态预测能力的问题，本文提出了基于自回归建模机制的人体姿态重构与预测方法 AIPose。该方法以当前 IMU 序列为输入，在统一框架下同时完成当前姿态重构、未来 IMU 序列预测和未来姿态预测。实验结果表明，在 XRFV2 数据集上，AIPose 当前姿态重构的 MPJPE 为 52.20\,px，PCK@5px、PCK@10px 和 PCK@20px 分别达到 16.13\%、26.86\% 和 44.55\%；未来姿态预测的 MPJPE 为 57.36\,px，PCK@5px、PCK@10px 和 PCK@20px 分别达到 16.48\%、26.65\% 和 43.51\%。DIP-IMU、AMASS 和 IMUPoser 数据集上的补充实验表明，AIPose 在不同 IMU 姿态数据上能够保持相对稳定的姿态建模能力；长时预测与扰动鲁棒性实验进一步说明，AIPose 在自回归预测过程中虽存在误差传播现象，但整体误差演化相对平稳。

总体来看，本文从人体活动识别与人体姿态重构与预测两个层面对非视觉人体感知问题进行了系统研究。第三章验证了 IMU 与 WiFi CSI 多模态对齐与融合对于提升活动识别性能的有效性，第四章则将研究重点从离散活动语义识别扩展到连续人体姿态建模，验证了基于 IMU 时序信号同时实现当前姿态重构与未来姿态预测的可行性。两部分研究在任务层次上相互衔接，共同体现了非视觉传感信号在隐私敏感、光照受限或视觉方法不适用场景中的应用潜力。

\xsection{展望}{Future Work}

尽管本文围绕人体活动识别与人体姿态重构与预测任务开展了较为系统的研究，并通过实验验证了所提方法的有效性，但从实验结果和实际应用需求来看，仍存在进一步改进空间。特别是在人体姿态重构与预测任务中，当前姿态重构精度、多任务优化平衡、跨场景与跨受试者泛化能力，以及长时预测稳定性和抗干扰能力仍需进一步提升。

首先，当前姿态重构性能与多任务优化平衡仍有改进空间。AIPose 需要同时优化当前姿态重构、未来 IMU 序列预测和未来姿态预测三个目标，因此在 DIP-IMU、AMASS 和 IMUPoser 等补充数据集上，与 ASIP、TIP 等专门面向当前姿态估计任务设计的方法相比，部分当前姿态重构指标仍存在差距。后续研究可从更强的 IMU 特征提取器、更精细的人体结构先验、更合理的姿态解码模块以及动态损失权重分配等方面进行改进，以减少多任务目标之间的相互干扰，并进一步提升整体姿态建模性能。

其次，人体姿态重构与预测模型的泛化能力仍需增强。第四章中的跨受试者与跨场景实验表明，当测试数据来自未参与训练的受试者或未见场景时，模型性能会出现明显下降，说明当前模型对个体运动习惯、身体结构差异、设备佩戴方式和环境分布变化仍较为敏感。未来可进一步探索面向 IMU 姿态建模的域泛化、少样本适配和测试时自适应方法，并结合自监督预训练或跨数据集联合训练，提升模型对不同数据来源、不同传感器配置和不同佩戴条件的适应能力。

最后，长时预测能力和抗干扰能力仍有待提高。本文实验表明，AIPose 在短时预测范围内能够保持相对稳定的性能，但随着预测时间延长，误差仍会逐渐累积。在真实应用中，传感器噪声、丢包、佩戴松动和人体运动突变等因素也可能进一步加剧误差传播。未来可引入层次化时序建模结构、多尺度预测机制、人体运动学约束或生成式预测模型，并在训练阶段加入噪声增强、传感器缺失模拟和扰动一致性约束，以提高模型在长时预测和复杂输入条件下的稳定性。

综上，未来工作可围绕当前姿态重构精度提升与多任务优化平衡、跨受试者和跨场景泛化能力增强、长时预测与抗干扰能力提升三个方向继续展开，从而推动非视觉人体感知方法向更高精度、更强泛化能力和更高实际可用性的方向发展。